\documentclass[10pt,a4paper]{amsart}
\usepackage[margin=19mm]{geometry}
\usepackage{amsmath,amssymb,mathtools}
\usepackage[scr=rsfs,cal=euler]{mathalfa}
\usepackage{xfrac}
\usepackage[unicode,hidelinks,bookmarks=false]{hyperref}
\newcommand{\ie}{i.e.\ }
\newcommand{\cf}{cf.\ }
\newcommand{\resp}{respectively\ }
\newcommand{\Subsection}{\subsection}
\providecommand{\nobreakdash}{\nobreak\hbox{-}}
\setlength{\emergencystretch}{2em}
\multlinegap0pt
\usepackage{mathtools}
\usepackage{amssymb}
\usepackage{float}
\usepackage{bm}
\usepackage{csquotes}
\usepackage{xcolor}
\newtheorem{theorem}{Theorem}
\newtheorem{corollary}{Corollary}
\newcommand{\D}{\mathcal{D}}
\newcommand{\E}{\mathbb{E}}
\newcommand{\F}{\mathcal{F}}
\newcommand{\clip}{\textnormal{clip}}
\renewcommand{\div}{\textnormal{div}}
\title{Posterior Concentration of Bayesian Physics-Informed Neural Networks for Elliptic PDEs}
\author{Yuxuan Zhao, Yulong Lu}
\date{Source: arXiv:2605.08672v1 (CC BY 4.0)}
\begin{document}
\maketitle

\noindent\emph{Source and licence.} Posterior Concentration of Bayesian Physics-Informed Neural Networks for Elliptic PDEs. Version arXiv:2605.08672v1 (CC BY 4.0), distributed by arXiv under the Creative Commons Attribution 4.0 International licence, http://creativecommons.org/licenses/by/4.0/, as recorded in the arXiv licence metadata for this exact version and shown on the abstract page https://arxiv.org/abs/2605.08672v1. Full licence notice: \texttt{licenses/arxiv-2605.08672-CC-BY-4.0.txt}.\par\smallskip

\noindent\textit{Editorial scope.} Counted unit: the article's corollary upper\_bound\_posterior\_mean (source0.tex lines 376--379). The article's complete original proof of it is delivered with it (source0.tex lines 1185--1253, one \texttt{proof} environment of 69 lines, which the article places away from the statement). No mathematics is written, repaired or re-derived: the statement and the proof are the article's own lines, in the article's own notation, and no retained line is substituted or re-flowed.  Retained uncounted as the source-local context the proof needs: lines 256--261; lines 336--338; lines 367--374; lines 610--612; lines 614--619; lines 621--626; lines 628--634; lines 1167--1173.  Nothing was omitted from any retained span, so the package carries no omission record.  No retained line contains a citation, so the wrapper carries no bibliography and no bibliography entry.  No retained line contains a figure environment, an image-inclusion command or drawing code, so no image is delivered and the package declares no asset dependency.  Editorial formatting only, in addition: an AMS \texttt{amsart} preamble that restates the article's own declarations and macros used by the retained lines, namely the environments \texttt{theorem}, \texttt{corollary} on separate counters, together with the standard packages those lines need.  The retained environments print in this document as Theorem 1, Corollary 1 in order of appearance, whereas the article prints the same items with the numbering of its own full document.  The complete native source of the article as distributed in the arXiv e-print for this exact version is bundled unchanged as \texttt{sources/200166/source0.tex}.
\begin{equation}\label{elliptic_pde}
\begin{aligned}
    - \div (A \nabla u) + V u = f  \quad & \textnormal{in} \ \Omega, \\
    u = g \quad & \textnormal{on} \ \partial \Omega, 
\end{aligned}
\end{equation}

\begin{align} \label{prior}
    \Pi(E) \coloneqq \Pi_\theta(\{\theta \in \Theta: \clip \circ f_\theta \in E \}), \quad E \in \Sigma_{\F}
\end{align}

\begin{theorem}\label{upper_bound}
     Assume that $u^*  \in C^\beta(\bar{\Omega})$ is the solution to $\eqref{elliptic_pde}$ with $f \in C^{\beta-2}(\bar{\Omega})$ and $g \in C^\beta(\partial{\Omega})$.  Assume that the smoothness level $\beta$ satisfies that  $2 < \beta \leq \beta^\ast$ for some $\beta^\ast > 2$, and $\|u^*\|_{C^{\beta}(\bar{\Omega})} \leq K$ for some $K > 0$. Then there exists 
     a prior $\Pi$ of the form $\eqref{prior}$ on a DNN-space $\F$ such that
     \begin{align} \label{posteior_contraction}
         \Pi(u \in \F: \mathcal{E}(u) > M_n^2 \epsilon_n^2 | \D^{(n)}) \to 0,
     \end{align}
     in $P_{u^*}^{(n)} Q_{u^*}^{(n)}$-probability as $n \to \infty$ for any $M_n \to \infty$. 
\end{theorem}

    \begin{align} \label{phi1}
        P_{u^*}^{(n)} \phi_1 \leq e^{C n \epsilon_n^2} \frac{e^{-\frac{1}{32}n M^2 \epsilon_n^2}}{1 - e^{-\frac{1}{32}n M^2 \epsilon_n^2}}, \quad P_u^{(n)} (1-\phi_1) \leq e^{-\frac{1}{32}n M^2 \epsilon_n^2 j^2},
    \end{align}

    \begin{equation} \label{conditioned_on_phi1}
    \begin{aligned}
        & P_{u^*}^{(n)} Q_{u^*}^{(n)} \left[\Pi(u \in \F_n: \|-\div (A \nabla (u - u^*)) + V(u - u^*) \|_{n, 1} >  J M \epsilon_n | \D^{(n)} ) \phi_1 \right] \\
        \leq & P_{u^*}^{(n)} Q_{u^*}^{(n)} ( \phi_1) = P_{u^*}^{(n)} (\phi_1) \leq e^{C n \epsilon_n^2} \frac{e^{-\frac{1}{32}n M^2 \epsilon_n^2}}{1 - e^{-\frac{1}{32}n M^2 \epsilon_n^2}}.
    \end{aligned}
    \end{equation}

    \begin{equation} \label{upper_bound_numerator}
        \begin{aligned}
            & P_{u^*}^{(n)} Q_{u^*}^{(n)} \left[ \int_{\|-\div (A \nabla (u - u^*)) + V(u - u^*) \|_{n, 1} \geq J M \epsilon_{n}} \frac{p_{u}^{(n)}}{p_{u^*}^{(n)}} \frac{q_{u}^{(n)}}{q_{u^*}^{(n)}} d \Pi(u)  (1 - \phi_{1}) \right] \\
     = & \int_{\|-\div (A \nabla (u - u^*)) + V(u - u^*) \|_{n, 1} \geq J M \epsilon_{n}}    P_{u}^{(n)} Q_{u}^{(n)} \left[(1 - \phi_{1}) \right] d \Pi(u) \leq  e^{-\frac{1}{36} n M^2 \epsilon_{n}^2 J^2}. 
        \end{aligned}
    \end{equation}

\begin{equation} \label{lower_bound_denominator}
    \begin{aligned}
        \int \frac{p_{u}^{(n)}}{p_{u^*}^{(n)}} \frac{q_{u}^{(n)}}{q_{u^*}^{(n)}} d \Pi(u) & > \int_{\mathcal{E}_n(u; \D^{(n)}) \leq \epsilon_n^2} \frac{p_{u}^{(n)}}{p_{u^*}^{(n)}} \frac{q_{u}^{(n)}}{q_{u^*}^{(n)}} d \Pi(u) \\
    & > \Pi(u \in \F: \mathcal{E}_n(u; \D^{(n)}) \leq \epsilon_n^2) e^{-(1+C_1) n\epsilon_n^2 } \\
    & > e^{-\xi n\epsilon_n^2 - (1+C_1) n \epsilon_n^2},
    \end{aligned}
\end{equation}

\begin{equation} \label{first_term_pupulation}
    \begin{aligned}
        & P_{u^*}^{(n)} Q_{u^*}^{(n)} \left[\Pi(u \in \F_n:  \|-\div(A \nabla (u-u^*)) + V(u-u^*) \|_{L^2(\Omega)} > 2 J M \epsilon_n| \D^{(n)})\right] \\
        \leq & P_{u^*}^{(n)} Q_{u^*}^{(n)} \left[\Pi(u \in \F_n: \|-\div(A \nabla (u-u^*)) + V(u-u^*) \|_{n, 1} >  J M \epsilon_n | \D^{(n)} ) \right] + e^{-n\epsilon_n^2} \\
        \leq & e^{C n \epsilon_n^2} \frac{e^{-\frac{1}{32}n M^2 \epsilon_n^2}}{1 - e^{-\frac{1}{32}n M^2 \epsilon_n^2}} + e^{-\frac{1}{36} n M^2 \epsilon_{n}^2 J^2 + \xi n\epsilon_n^2 + (1+C_1) n \epsilon_n^2} + C_1^{-1} (n\epsilon_n^2)^{-1} + e^{-n\epsilon_n^2}. 
    \end{aligned}
\end{equation}

\begin{corollary} \label{upper_bound_posterior_mean}
    Assume the same setting as in Theorem \ref{upper_bound}. Moreover, let $\Pi$ and the DNN-space $\F$ to be the same as in Theorem \ref{upper_bound}. Define the posterior mean as $\bar{u}_n \coloneqq \E_{u \in \Pi} \left[u | \D^{(n)}\right]$. Then with high probability,  $\mathcal{E}(\bar{u}_n)  \lesssim \epsilon_n^2$ as $n \rightarrow \infty$. 
    
\end{corollary}

\begin{proof}[Proof of Corollary \ref{upper_bound_posterior_mean}]
    Note that it suffices to show that for any $M_n \to \infty$,
    \begin{align*}
        P_{u^*}^{(n)} Q_{u^*}^{(n)} \left(u \in \F: \|-\div(A \nabla \bar{u}_n ) + V \bar{u}_n  -f\|_{L^2(\Omega)}^2 + \lambda \|\bar{u}_n - g\|_{L^2(\partial \Omega)}^2 \geq M_n^2 \epsilon_n^2\right) \to 0,
    \end{align*}
    as $n \to \infty$. 
    Also by Jensen's inequality, we have
    \begin{align*}
         \|-\div(A \nabla \bar{u}_n ) + V \bar{u}_n  -f\|_{L^2(\Omega)} \leq \E_{u \sim \Pi} \left[ \|-\div(A \nabla u) + V u -f\|_{L^2(\Omega)} | \D^{(n)}\right], 
    \end{align*}
    and 
    \begin{align*}
        \|\bar{u}_n - g\|_{L^2(\partial \Omega)} \leq \E_{u \sim \Pi} \left[\|u -g\|_{L^2(\partial \Omega)} | \D^{(n)}\right]. 
    \end{align*}
    Hence to prove the statement, it is sufficient to show that for any $M_n \to \infty$
    \begin{equation} \label{posterior_mean_domain}
         P_{u^*}^{(n)} Q_{u^*}^{(n)} \left( \E_{u \sim \Pi} \left[ \|-\div(A \nabla u) + V u -f\|_{L^2(\Omega)} | \D^{(n)}\right] > M_n \epsilon_n \right) \to 0,
    \end{equation}
    and 
    \begin{equation} \label{posterior_mean_boundary}
        P_{u^*}^{(n)} Q_{u^*}^{(n)} \left(\E_{u \sim \Pi} \left[\|u -g\|_{L^2(\partial \Omega)} | \D^{(n)}\right] > M_n \epsilon_n \right) \to 0,
    \end{equation}
    as $n \to \infty$.
    
    By Cauchy-Schwarz and Jensen's inequality, 
    \begin{equation} \label{posterior_mean_upper_bound}
        \begin{aligned}
            & \E_{u \sim \Pi} \left[  \|-\div(A \nabla u) + V u -f\|_{L^2(\Omega)} | \D^{(n)} \right] \\ \leq &  2J M \epsilon_n  + \E_{u \sim \Pi}\left[  \|-\div(A \nabla u) + V u -f\|_{L^2(\Omega)} \mathbf{1}_{\{  \|-\div(A \nabla u) + V u -f\|_{L^2(\Omega)} >2 J M \epsilon_n \}} | \D^{(n)}\right] \\
              \leq & 2J M \epsilon_n + \left(\E \left[ \|-\div(A \nabla u) + V u -f\|_{L^2(\Omega)}^2 | \D^{(n)}\right] \right)^{1/2} \left(\Pi \left[ \|-\div(A \nabla u) + V u -f\|_{L^2(\Omega)} >2 J M \epsilon_n | \D^{(n)} \right]\right)^{1/2}. 
        \end{aligned}
    \end{equation}
    Let $\phi_1$ be test on $\{X_i, f_i\}_{i=1}^n$ such that \eqref{phi1} holds, and let $E$ to be the set on which \eqref{lower_bound_denominator} holds. Recall that by \eqref{conditioned_on_phi1}, \eqref{upper_bound_numerator} and \eqref{lower_bound_denominator} 
    \begin{equation*}
        \begin{aligned}
            & P_{u^*}^{(n)} Q_{u^*}^{(n)} \left[\Pi(u \in \F_n: \|-\div(A \nabla (u-u^*)) + V (u-u^*) \|_{n, 1} >  J M \epsilon_n | \D^{(n)} ) \mathbf{1}_E \right] \\
            \leq &   P_{u^*}^{(n)} Q_{u^*}^{(n)} (\phi_1) + P_{u^*}^{(n)} Q_{u^*}^{(n)} \left[\Pi(u \in \F_n: \|-\div(A \nabla (u-u^*)) + V (u-u^*)\|_{n, 1} >  J M \epsilon_n | \D^{(n)} ) (1 - \phi_1)\mathbf{1}_E \right] \\
            \leq & e^{C n \epsilon_n^2} \frac{e^{-\frac{1}{32}n M^2 \epsilon_n^2}}{1 - e^{-\frac{1}{32}n M^2 \epsilon_n^2}} + e^{-\frac{1}{36} n M^2 \epsilon_{n}^2 J^2 + \xi n\epsilon_n^2 + (1+C_1) n \epsilon_n^2}. 
        \end{aligned}
    \end{equation*}
    Then by \eqref{first_term_pupulation}, we can take $J, M$ to be large enough such that 
    \begin{equation}
    \begin{aligned}
        & P_{u^*}^{(n)} Q_{u^*}^{(n)} \left[\Pi(u \in \F_n: \|-\div(A \nabla (u-u^*)) + V (u-u^*) \|_{L^2(\Omega)} > 2 J M \epsilon_n| \D^{(n)}) \mathbf{1}_E \right] \\
        \leq & P_{u^*}^{(n)} Q_{u^*}^{(n)} \left[\Pi(u \in \F_n: \|-\div(A \nabla (u-u^*)) + V (u-u^*) \|_{n, 1} >  J M \epsilon_n | \D^{(n)} ) \mathbf{1}_E \right] + e^{-  n\epsilon_n^2 J^2} \\
        \leq & e^{C n \epsilon_n^2} \frac{e^{-\frac{1}{32}n M^2 \epsilon_n^2}}{1 - e^{-\frac{1}{32}n M^2 \epsilon_n^2}} + e^{-\frac{1}{36} n M^2 \epsilon_{n}^2 J^2 + \xi n\epsilon_n^2 + (1+C_1) n \epsilon_n^2} + e^{- n\epsilon_n^2 J^2}. 
    \end{aligned}
\end{equation}
Furthermore, by Markov's inequality 
\begin{equation} \label{markov_upper_bound}
    \begin{aligned}
        & P_{u^*}^{(n)} Q_{u^*}^{(n)} \left(\Pi(u \in \F_n: \|-\div(A \nabla u) + V u -f\|_{L^2(\Omega)} > 2 J M \epsilon_n| \D^{(n)}) \mathbf{1}_E > e^{-b n\epsilon_n^2} \right) \\
        \leq & e^{b n\epsilon_n^2} \left(e^{C n \epsilon_n^2} \frac{e^{-\frac{1}{32}n M^2 \epsilon_n^2}}{1 - e^{-\frac{1}{32}n M^2 \epsilon_n^2}} + e^{-\frac{1}{36} n M^2 \epsilon_{n}^2 J^2 + \xi n\epsilon_n^2 + (1+C_1) n \epsilon_n^2} + e^{-  n\epsilon_n^2 J^2}\right). 
    \end{aligned}
\end{equation}
Now we can estimate the second term in \eqref{posterior_mean_upper_bound} by 
\begin{equation} \label{probability_second_term}
    \begin{aligned}
        & P_{u^*}^{(n)} Q_{u^*}^{(n)} \left(\E \left[\|-\div(A \nabla u) + V u -f\|_{L^2(\Omega)}^2 | \D^{(n)}\right] \cdot \Pi \left[ \|-\div(A \nabla u) + V u -f\|_{L^2(\Omega)} > 2 J M \epsilon_n | \D^{(n)} \right] > \epsilon_n^2 \right) \\
        \leq &  P_{u^*}^{(n)} Q_{u^*}^{(n)} \left(\E \left[ \|-\div(A \nabla u) + V u -f\|_{L^2(\Omega)}^2 | \D^{(n)}\right] \cdot \Pi \left[ \|-\div(A \nabla u) + V u -f\|_{L^2(\Omega)} >2 J M \epsilon_n | \D^{(n)} \right] \mathbf{1}_E > \epsilon_n^2 \right) \\ & + P_{u^*}^{(n)} Q_{u^*}^{(n)} (\mathbf{1}_{E^c}) \\
        \leq & P_{u^*}^{(n)} Q_{u^*}^{(n)} \left(\E \left[ \|-\div(A \nabla u) + V u -f\|_{L^2(\Omega)}^2 | \D^{(n)}\right] \mathbf{1}_E  \cdot e^{-b n\epsilon_n^2 } > \epsilon_n^2  \right) \\
        & + e^{b n\epsilon_n^2} \left(e^{C n \epsilon_n^2} \frac{e^{-\frac{1}{32}n M^2 \epsilon_n^2}}{1 - e^{-\frac{1}{32}n M^2 \epsilon_n^2}} + e^{-\frac{1}{36} n M^2 \epsilon_{n}^2 J^2 + \xi n\epsilon_n^2 + (1+C_1) n \epsilon_n^2} + e^{- n\epsilon_n^2 J^2}\right) + C_1^{-1} (n \epsilon_n^2)^{-1}\\
        = &  P_{u^*}^{(n)} Q_{u^*}^{(n)} \left( \frac{\int \|-\div(A \nabla u) + V u -f\|_{L^2(\Omega)}^2 \frac{p_{u}^{(n)}}{p_{u^*}^{(n)}} \frac{q_{u}^{(n)}}{q_{u^*}^{(n)}} d \Pi(u)}{\int \frac{p_{u}^{(n)}}{p_{u^*}^{(n)}} \frac{q_{u}^{(n)}}{q_{u^*}^{(n)}} d \Pi(u)} \mathbf{1}_E > e^{b n \epsilon_n^2} \epsilon_n^2 \right) \\
        & + e^{b n\epsilon_n^2} \left(e^{C n \epsilon_n^2} \frac{e^{-\frac{1}{32}n M^2 \epsilon_n^2}}{1 - e^{-\frac{1}{32}n M^2 \epsilon_n^2}} + e^{-\frac{1}{36} n M^2 \epsilon_{n}^2 J^2 + \xi n\epsilon_n^2 + (1+C_1) n \epsilon_n^2} + e^{-  n\epsilon_n^2 J^2}\right) + C_1^{-1} (n \epsilon_n^2)^{-1} \\
        \leq & e^{\xi n\epsilon_n^2 + (1+C_1) n \epsilon_n^2} e^{-b n\epsilon_n^2} \epsilon_n^{-2} \int \|-\div(A \nabla u) + V u -f\|_{L^2(\Omega)}^2 d \Pi(u) \\
        & + e^{b n\epsilon_n^2} \left(e^{C n \epsilon_n^2} \frac{e^{-\frac{1}{32}n M^2 \epsilon_n^2}}{1 - e^{-\frac{1}{32}n M^2 \epsilon_n^2}} + e^{-\frac{1}{36} n M^2 \epsilon_{n}^2 J^2 + \xi n\epsilon_n^2 + (1+C_1) n \epsilon_n^2} + e^{-  n\epsilon_n^2 J^2}\right) + C_1^{-1} (n \epsilon_n^2)^{-1}. 
    \end{aligned}
\end{equation}
where the second inequality follows from \eqref{markov_upper_bound}, the third inequality follows from Markov's inequality, Fubini's theorem and \eqref{lower_bound_denominator}. Then we can take $C_1 = 1$, $b = \xi + 1 + C_1$ and $M, J$ to be some large constant such that  \eqref{probability_second_term} goes to $0$ as $n \to \infty$. Together with \eqref{posterior_mean_upper_bound}, we prove $\eqref{posterior_mean_domain}$. \eqref{posterior_mean_boundary} can argued in a similar way. Therefore we conclude the proof. 
\end{proof}

\end{document}
